\section{The Monk at the Helm of a Dying Rome}

Eight months after his consecration as bishop of Rome, Gregory wrote to his
closest friend in Spain, Bishop Leander of Seville, to describe what the papacy
felt like from within. He did not write of triumph or majesty. He wrote of a
failing hull in heavy seas:

\begin{quote}
For, indeed, I am in this place tossed by such billows of this world that I am
in no wise able to steer into port the old and rotten ship of which, in the
hidden dispensation of God, I have assumed the guidance. Now in front the
billows rush in, now at the side heaps of foamy sea swell up, now from behind
the storm follows on. \dots\ I groan, because I feel that through my negligence
the bilgewater of vices increases, and, as the storm meets the vessel
violently, the rotten planks already sound of shipwreck. With tears I remember
how I have lost the placid shore of my rest, and with sighs I behold the land
which still, with the winds of affairs blowing against me, I cannot reach.
(\work{Ep.}\ I.43, May 591)
\end{quote}

Every enduring question about Pope Saint Gregory the Great begins inside that
image. Born around 540 into a Rome shattered by Gothic war, depopulation, and
famine, he surrendered the highest civic office in the city to become a monk on
his family's estate on the Caelian Hill. Yet monastic silence never kept him.
Pulled back into the world first as papal envoy at the imperial court of
Constantinople and then, in the plague year 590, forced onto the chair of Saint
Peter, he spent fourteen years governing a besieged city, managing estates from
Sicily to Gaul to feed refugees and ransom captives, defending episcopal
humility against Constantinopolitan titles, dispatching monks to the English,
and preaching Scripture while Lombard horsemen rode beneath the Roman walls.

Between the sixth-century bishop and the figure honored by later Christendom
stands a vast curtain of memory. Within two centuries of his death, England
remembered him as its apostle through the famous wordplay of the Roman slave
market; Carolingian Rome and Francia attached his name to the papal
sacramentary and the full repertoire of Latin plainchant; medieval legend
painted a white dove whispering at his ear, set the Archangel Michael sheathing
a fiery sword above Hadrian's mausoleum, and told of his tears winning mercy
for the pagan Emperor Trajan. Those traditions reveal how deeply Western
Christendom loved him, yet they easily crowd out the far more compelling man
preserved in his own letters, homilies, and treatises and in the eyewitness
report of his contemporary Gregory of Tours.

How did a contemplative Roman aristocrat who believed the world was growing old
and nearing its end become the most effective pastor, administrator, and
missionary strategist of early medieval Europe? What did his own writings and
contemporaries actually witness, and where did his character show genuine
strain---in his dependence on imperial coercion against schismatics, his
ignorance of Greek, and his troubling letters to the usurper Phocas? And how
did the real historical work of 590--604 grow into the monumental legacy of the
Latin Doctor whose name still marks the Mass, the chant, and the pastoral care
of the Church?
